% !TeX root = ../nc-del-pezzo.tex
\section{次数1の全係数Hessian候補と特殊族からの決定の限界}
\label{sec:kuznetsov-hessian-counterexample}

本節では，全係数に対する二次差の候補を実際のカテゴリー的二次差から区別する。
巡回族のHessian公式は，自然な全係数多項式へ拡張できる。
しかし，その候補には滑らかな曲面上で零となる非零方向が存在し，
さらに全ての既知の巡回族計算を保つ別の多項式項も存在する。
従って，一般公式の確定には新しいカテゴリー的な計算が必要である。

\subsection{全係数に対する自然な二次共変式}

次数1の曲面を
\begin{equation}\label{eq:ku-d1-full-hessian-model}
 X:\quad F=\frac{z^2}{2}+f(x,u,y)=0,
 \qquad f=\frac{y^3}{3}+A(x,u)y+B(x,u),
 \quad \deg A=4,\quad\deg B=6
\end{equation}
と書く。重みは $(1,1,2,3)$ であり，$X$ は滑らかとする。
$D=d_x\partial_x+d_u\partial_u$ をPoisson方向の表示とする。
重み6のJacobian商による古典変形の表示では
\begin{equation}\label{eq:ku-d1-reduced-jacobian}
 \begin{split}
 J_{F,6}
 &:=\left(\frac{\C[x,u,y,z]}{(F_x,F_u,F_y,F_z)}\right)_6\\
 &\simeq
 \frac{\Sym^6\C^2{}^*\oplus y\Sym^4\C^2{}^*}
 {\{V(B)+yV(A):V\in\mathfrak{gl}_2\}}.
 \end{split}
\end{equation}
実際，$F_z=z$ と $F_y=y^2+A$ を先に消す。
残る重み5の二つの関係の重み6部分は，それぞれに一次式を掛けたものなので，
分母は表示した線型ベクトル場の像に一致する。

\begin{proposition}[次数1のHessian共変式]
\label{prop:ku-d1-full-hessian-covariant}
次の式は，古典接空間に値を持つ自然な二次テンソルを定める：
\begin{equation}\label{eq:ku-d1-full-hessian-covariant}
 \mathcal H_{1,X}(D)
 :=\frac18\left[2y\bigl(yD^2A+D^2B\bigr)-(DA)^2\right]
 =\frac18\left[2yD^2B-2A D^2A-(DA)^2\right].
\end{equation}
特に $D=\partial_x$ なら，最初の代表元は
$\det\operatorname{Hess}_{x,y}(f)/8$ である。
巡回族 $A=\alpha x^4$，$B=\beta x^6+u^6$，$D=\partial_x$ では
\[
 \mathcal H_{1,X}(D)
 =\left[\frac{15}{2}\beta x^4y-5\alpha^2x^6\right]
 =\delta_{2,X}(D,D).
\]
最後の等式は定理\ref{thm:ku-degree-one-quadratic-obstruction}による。
一般の $(A,B,D)$ に対する $\mathcal H_{1,X}(D)=\delta_{2,X}(D,D)$ は
本命題の主張に含めない。
\end{proposition}
\begin{proof}
二つの表示は $y^2=-A$ により一致する。
depressed表示の座標変換を
\[
 \binom{x'}{u'}=M\binom{x}{u},\quad y'=cy,\quad z'=\eta z,
 \qquad \eta^2=c^3
\]
とする。このとき
\[
 A'=\frac{\eta^2}{c}A\circ M^{-1},\qquad
 B'=\eta^2B\circ M^{-1},\qquad
 D'=\frac c\eta MD.
\]
連鎖律から，変数の引戻しを略記して
\[
 D'A'=\eta DA,\qquad (D')^2A'=cD^2A,
 \qquad (D')^2B'=c^2D^2B
\]
となる。従って式\eqref{eq:ku-d1-full-hessian-covariant}の代表元は
方程式の変形と同じ $\eta^2$ 倍の変換則を持ち，Jacobian商へ降下する。
巡回族では $D^2A=12\alpha x^2$，$DA=4\alpha x^3$，
$D^2B=30\beta x^4$ を代入すればよい。
\end{proof}

\subsection{滑らかな曲面上での非零方向の消滅}

\begin{proposition}[Hessian候補の零方向を持つ滑らかな曲面]
\label{prop:ku-d1-hessian-null-direction}
式\eqref{eq:ku-d1-full-hessian-model}において
\begin{equation}\label{eq:ku-d1-hessian-null-example}
 \begin{split}
 A={}&x^3u-\frac{21}{2}x^2u^2+21xu^3-14u^4,\\
 B={}&-\frac1{20}x^5u+\frac72x^4u^2-\frac{49}{2}x^3u^3
       +70x^2u^4-\frac{161}{2}xu^5+\frac{343}{10}u^6
 \end{split}
\end{equation}
と置く。この $X$ は $\mathbb Q$ 上で定義された滑らかな次数1 del Pezzo曲面であり，
非零Poisson方向 $D=\partial_x$ に対して
\[
                         \mathcal H_{1,X}(D)=0
\]
となる。
\end{proposition}
\begin{proof}
$V=x\partial_x-5u\partial_u$ と置く。
式\eqref{eq:ku-d1-hessian-null-example}の直接微分により
\begin{equation}\label{eq:ku-d1-hessian-null-identities}
 2B_{xx}=V(A),\qquad -2AA_{xx}-A_x^2=V(B)
\end{equation}
を得る。従って多項式環の中で既に
\begin{equation}\label{eq:ku-d1-hessian-exact-syzygy}
 \det\operatorname{Hess}_{x,y}(f)
 =V(f)+2A_{xx}f_y
 =xf_x-5uf_u+2A_{xx}f_y
\end{equation}
である。これはJacobianイデアルの元なので，候補の消滅が従う。

滑らかさも厳密に確認する。
$u=1$ で $\Delta(x)=4A(x,1)^3+9B(x,1)^2$ と置くと，
$400\Delta(x)=P(x)$ であり
\begin{equation}\label{eq:ku-d1-hessian-discriminant}
 \begin{split}
 P(x)={}&9x^{10}+340x^9+2520x^8-12600x^7-82320x^6\\
       &+252252x^5+246960x^4-1228920x^3\\
       &+1102500x^2-123480x-155036.
 \end{split}
\end{equation}
この多項式は $\mathbb Q$ 上平方因子を持たない。
例えば $\mathbf F_{11}[x]$ で，次の二つの多項式が
$U\overline P+W\overline P'=1$ を満たす：
\begin{equation}\label{eq:ku-d1-hessian-bezout}
 \begin{split}
 U={}&-5x^8+5x^7+5x^6+x^5-5x^4+4x^3+4x^2+5x-3,\\
 W={}&-5x^9+2x^8-5x^7-3x^6-4x^5-4x^4
           +5x^3-4x^2+1.
 \end{split}
\end{equation}
分母60と $P$ の最高次係数9は11で可逆なので，この証明書は
$\mathbb Q$ 上の平方因子の不存在を証明する。

$u=1$ 上に特異点があれば，$F_z=0$ から $z=0$ であり，
$f=f_y=f_x=0$ となる。
その点で $A=-y^2$，$B=2y^3/3$，$B_x=-yA_x$ だから
\[
 \Delta=0,\qquad
 \Delta_x=12A^2A_x+18BB_x=0
\]
となり，平方因子の不存在に矛盾する。
他方，$u=0$ では $A=B=0$ である。
全ての偏微分が零になる可能性がある点は $y=z=0$，$x\ne0$ だけだが，
そこでは
\[
                             F_u=-x^5/20\ne0
\]
である。$x=u=0$ には錐の臨界点はない。
また $y^3,z^2$ の係数は非零なので，$X$ は重み付き射影空間の
二つの特異点を通らない。従って $X$ は滑らかである。
重み付き随伴公式から $-K_X=\OO_X(1)$，$K_X^2=1$ を得る。
\end{proof}

\begin{remark}[候補の反例と実際の二次差]
\label{rem:ku-d1-hessian-null-scope}
命題\ref{prop:ku-d1-hessian-null-direction}は
全係数Hessian候補に関する証明済みの反例である。
この曲面で実際の $\delta_{2,X}(D,D)$ が零になることや，
Hessian候補が実際の二次差と一致しないことは，まだ結論しない。
また，系\ref{cor:ku-generic-all-directions-quadratic}の
稠密開上の全非零方向での非消滅とは矛盾しない。
一般公式がHessian候補に一致する場合には，その稠密開から除くべき
滑らかな特殊点が実際に存在することになる。
\end{remark}

\subsection{巡回族では検出できない多項式項}

\begin{proposition}[自然性と多項式性だけでは次数1の式は決まらない]
\label{prop:ku-d1-polynomial-covariant-ambiguity}
\begin{equation}\label{eq:ku-d1-hidden-covariant}
 \mathcal R_{1,X}(D):=\left[4A D^2A-3(DA)^2\right]\in J_{F,6}
\end{equation}
は自然な多項式二次テンソルである。
さらに $\ell_D=d_xu-d_ux$ と置くと，多項式環の中で
\begin{equation}\label{eq:ku-d1-hidden-covariant-binary-hessian}
 4A D^2A-3(DA)^2
 =\frac13\ell_D^2\det\operatorname{Hess}_{x,u}(A)
\end{equation}
が成立する。
これは $A=\alpha\ell^4$ の全ての点と全ての方向で零となるが，
滑らかな次数1曲面のモジュライ上で恒等的には零でない。
従って各 $s\in\C$ に対する
\[
                  \mathcal H_{1,X}+s\mathcal R_{1,X}
\]
は互いに異なる自然な多項式候補であり，全ての既知の巡回族とFermat点の計算に一致する。
特に，滑らかな係数空間上の正則性を全係数空間上の多項式性に強めても，
これらの計算と座標変換則だけでは実際の二次差を一意に決められない。
\end{proposition}
\begin{proof}
命題\ref{prop:ku-d1-full-hessian-covariant}の変換則から
$A'(D')^2A'=\eta^2A D^2A$ と $(D'A')^2=\eta^2(DA)^2$ となるので，
式\eqref{eq:ku-d1-hidden-covariant}は古典接方向として自然である。
式\eqref{eq:ku-d1-hidden-covariant-binary-hessian}は $D=\partial_x$ の場合を
確認すれば，座標変換則から一般に従う。
Euler恒等式を微分すると
\[
 uA_{xu}=3A_x-xA_{xx},\qquad
 u^2A_{uu}=12A-6xA_x+x^2A_{xx}
\]
となり，これらを $u^2(A_{xx}A_{uu}-A_{xu}^2)$ へ代入すれば
$12AA_{xx}-9A_x^2$ を得る。
$A=\alpha\ell^4$ なら
\[
 4A D^2A=48\alpha^2\ell^6(D\ell)^2=3(DA)^2
\]
であり，Jacobian商を取る前から零になる。

非消滅の簡単な例は
\[
 X_0:\quad z^2/2+y^3/3+(x^4+u^4)y+x^6+u^6=0,
 \qquad D=\partial_x
\]
である。この曲面は滑らかである。実際，$F_z=0$ から $z=0$，
残る偏微分は
\[
 F_x=2x^3(2y+3x^2),\quad F_u=2u^3(2y+3u^2),
 \quad F_y=y^2+x^4+u^4
\]
である。$xu\ne0$ で最初の二つが零なら $x^2=u^2$，$y=-3x^2/2$ なので，
最後は $17x^4/4\ne0$ である。一方だけが非零なら最後は
$13x^4/4$ または $13u^4/4$ となり，双方が零の場合も臨界点はない。
この例で $\mathcal R_{1,X_0}(D)=48[x^2u^4]$ である。
式\eqref{eq:ku-d1-reduced-jacobian}により，純粋な六次式の類が零なら
$V(A)=0$，$V(B)$ がその六次式に等しい線型ベクトル場 $V$ が必要である。
$A=x^4+u^4$ については $V(A)=0$ の係数を比較すると $V=0$ となるので，
$48[x^2u^4]$ は非零である。

さらに命題\ref{prop:ku-d1-hessian-null-direction}の曲面そのものでも
$\mathcal R_{1,X}(\partial_x)\ne0$ となる。代表元は
\[
 -3u^2\bigl(x^4-14x^3u+105x^2u^2-182xu^3+49u^4\bigr)
\]
である。この $A$ に対して，線型ベクトル場の基底
$x\partial_x,u\partial_x,x\partial_u,u\partial_u$ と
係数 $x^4,x^3u,x^2u^2,xu^3$ を用いた $V\mapsto V(A)$ の小行列は
\[
 \begin{pmatrix}
 0&0&1&0\\
 3&0&-21&1\\
 -21&3&63&-21\\
 21&-21&-56&63
 \end{pmatrix},\qquad\det=-378.
\]
従ってここでも $V(A)=0$ は $V=0$ を強制する。
この同じ滑らかな点で，$s=0$ の候補だけが指定方向で零になる。
\end{proof}

式\eqref{eq:ku-d1-hidden-covariant-binary-hessian}から，
この追加項は中央一次式 $\ell_D=0$ に制限すると零になる。
従ってその境界上へ制限した方程式だけの比較では，
候補 $\mathcal H_1+s\mathcal R_1$ の定数 $s$ を区別できない。
境界の高い近傍や同変周期の全情報が同じになるという主張ではない。

\subsection{多項式延長を仮定した全候補の分類}

$W=\C^2$ とし，$A\in\Sym^4W^*$，$B\in\Sym^6W^*$ の
\emph{全係数空間}の座標環を $R$ とする。
ここでは特異な方程式も含める。
普遍Jacobian加群の重み6部分を
\[
 \mathscr J_6:=
 \left(\frac{R[x,u,y]}{(y^2+A,\ yA_x+B_x,\ yA_u+B_u)}\right)_6
\]
と書く。これは全係数空間では局所自由とは限らない。
$G=\operatorname{GL}(W)\times\Gm$ は上の座標変換により作用する。
後者の因子を $c=t^2,\eta=t^3$ と表示すると
\[
 A\mapsto t^4A,\quad B\mapsto t^6B,\quad
 y\mapsto t^2y,\quad D\mapsto t^{-1}D
\]
である。

\begin{proposition}[多項式延長の下で残る一つの定数]
\label{prop:ku-d1-polynomial-classification}
$(A,B)$ に多項式的に依存し，$D$ に斉次二次である
$G$-共変な $\mathscr J_6$ 値の式の空間は三次元であり，
\begin{equation}\label{eq:ku-d1-polynomial-basis}
                 [yD^2B],\qquad[A D^2A],\qquad[(DA)^2]
\end{equation}
を基底に持つ。
そのうち定理\ref{thm:ku-degree-one-quadratic-obstruction}の
全巡回族での値に一致するものは，ちょうど
\begin{equation}\label{eq:ku-d1-conditional-general-form}
                         \mathcal H_1+s\mathcal R_1,
                         \qquad s\in\C
\end{equation}
である。

特に，実際の $\delta_2$ がこの普遍Jacobian加群へ共変多項式として
延長するという\emph{追加仮定}の下では，その全係数公式は
式\eqref{eq:ku-d1-conditional-general-form}の一つに限られる。
この延長の存在と定数 $s$ の値は，いずれも本命題では証明していない。
\end{proposition}
\begin{proof}
$y^2=-A$ により，共変式は
$R\otimes(\Sym^6W^*\oplus y\Sym^4W^*)$ の商に値を持つ。
商における共変式は，この自由加群に値を持つ共変多項式へ持ち上がる。
実際，任意の多項式の持上げを取り，その $G$ 軌道が張る有限次元有理表現の中で
極大compact部分群に沿って平均すればよい。
平均は商に与えた共変式を変えず，compact部分群のZariski稠密性から
$G$ に対する共変性を持つ。
ここで多項式延長の仮定を使っている。

代表元の $y^e$ 係数，$e=0,1$ において，$A,B$ に対する多項式次数を
それぞれ $i,j$ とする。$D$ の次数が2であることと $t$ の重みを比較すると
\[
                       4i+6j-2=6-2e
\]
である。従って純六次式の部分は $A$ に二次であり，
$y$ の係数は $B$ に一次である。
対応する $\operatorname{GL}(W)$ のintertwiner空間は
\[
 \begin{split}
 &\Hom_{\operatorname{GL}(W)}
   \bigl(\Sym^2(\Sym^4W^*)\otimes\Sym^2W,\Sym^6W^*\bigr),\\
 &\Hom_{\operatorname{GL}(W)}
   \bigl(\Sym^6W^*\otimes\Sym^2W,\Sym^4W^*\bigr)
 \end{split}
\]
であり，次元はそれぞれ2と1である。
この小さい表現計算はweightだけでも確認できる。
有限次元 $\operatorname{SL}_2$ 表現において，最高weight $\ell$ の表現の重複度は
weight $\ell$ と $\ell+2$ の次元差である。
第一のsourceではweight6,8の次元が4,2，
第二のsourceではweight4,6の次元が3,2となる。
中心のweightもそれぞれ $-6,-4$ なので，これが上の
$\operatorname{GL}_2$ の重複度である。
第一の二つの独立な式は $A D^2A,(DA)^2$，第二は $D^2B$ である。
従って式\eqref{eq:ku-d1-polynomial-basis}は商の全共変式を張る。

商を取ってもこの三つは独立である。
まず $A=0,B=x^6+u^6,D=\partial_x$ では，
$[yD^2B]=30[x^4y]\ne0$ で，他の二つは零である。
次に $A=x^4,B=u^6,D=\partial_x$ では
$[A D^2A]=12[x^6]$，$[(DA)^2]=16[x^6]$，$[x^6]\ne0$ である。
これらの非消滅は式\eqref{eq:ku-d1-reduced-jacobian}から直接分かり，
二つの曲面はいずれも滑らかである。
従って定数係数の線型関係が存在すれば，それは
$4[A D^2A]-3[(DA)^2]=\mathcal R_1$ の倍数に限られる。
命題\ref{prop:ku-d1-polynomial-covariant-ambiguity}によりこの最後の式も非零なので，
三つは独立である。

最後に一般の候補を $p[yD^2B]+q[A D^2A]+r[(DA)^2]$ と書く。
同じ二つの巡回族の点で既知の二次差と比較すると
\[
                             p=\frac14,\qquad12q+16r=-5
\]
となる。この解は
$q=-1/4+4s,r=-1/8-3s$ であり，
式\eqref{eq:ku-d1-conditional-general-form}に一致する。
\end{proof}

以上の多項式恒等式，判別式，有限体B\'ezout証明書，小行列式は
\texttt{computations/verify\_degree\_one\_hessian\_counterexample.py}
で有理数・整数の厳密演算により再現できる。
この計算は候補の性質を検証するものであり，
非可換Hochschild--Serre乗法や実際の $\delta_2$ の一般公式を検証したものではない。
